\documentclass{article}
\usepackage{amsmath, amssymb, booktabs}

\begin{document}

\section*{PM vs GR: Cross-Force Interaction Comparison}

This document directly compares what the Pushing--Medium (PM) theory predicts
regarding interactions between gravitational and electromagnetic forces versus
what General Relativity (GR) can realize.

\section*{1. Side-by-Side Comparison}

\begin{center}
\begin{tabular}{p{4cm}|p{5cm}|p{5cm}}
\toprule
\textbf{Physical Scenario} & \textbf{General Relativity} & \textbf{Pushing-Medium} \\
\midrule
\textbf{Charged mass lensing} & 
Lensing from total mass-energy only. Charge contributes via $E^2/(2\varepsilon_0)$ energy—negligible. &
Lensing from $\delta n_G + \delta n_E$. For charge $q$, mass $M$: additive contributions with ratio $(\mu_E/\mu_G)(q/M)(G/k_e)$. \\
\midrule
\textbf{Light through plasma} & 
Plasma affects light speed via dispersion $n_{\text{plasma}}(\omega)$. Gravity and plasma are independent. &
Plasma modifies \emph{gravitational} lensing: $n_{\text{tot}} = 1 + \mu_G\Phi_G - \alpha_{\text{plasma}}n_e$. Gravitational deflection depends on electron density. \\
\midrule
\textbf{Moving charged mass} & 
Frame-dragging from total momentum $T^{0i}$. EM and gravitational contributions separate. &
Flow field $\mathbf u = \mathbf u_G + \mathbf u_q$ with contributions from both mass current $\mathbf J_M$ and charge current $\mathbf J_q$. \\
\midrule
\textbf{Electric current near light beam} & 
No gravitational effect beyond stress-energy ($\sim B^2 c / \mu_0$, utterly negligible). &
Creates flow field $\nabla^2\mathbf u_q = -A_q\mathbf J_q$. Drags light via Fizeau term. Constrained to $A_q \lesssim 10^{-10}$ m$^3$/(C$\cdot$s). \\
\midrule
\textbf{Magnetic field gradient} & 
Gravitates via field energy $B^2/(2\mu_0)$. No direct geometric effect. &
Initially suggested $\mathbf u_B = \kappa_B\mathbf A$, but \emph{ruled out} by lab null results. Sets $\kappa_B \approx 0$. \\
\midrule
\textbf{Dynamic coupling} & 
$\partial_t g_{\mu\nu}$ depends on $T_{\mu\nu}$. EM and gravity evolve independently except via energy-momentum. &
Advection $\nabla\cdot(\phi\mathbf u)$ couples scalar and vector fields dynamically. Since $\mathbf u$ includes all currents, energy transfers between force sectors. \\
\midrule
\textbf{Gravitational wave speeds} & 
Single speed $c$ for all gravitational radiation (tensor modes). &
Two speeds $c_\phi$ and $c_u$ for scalar and vector medium waves. Could lead to multi-component GW signals. \\
\bottomrule
\end{tabular}
\end{center}

\section*{2. Concrete Quantitative Examples}

\subsection*{2.1 Electron Lensing}

\textbf{Scenario:} An electron ($m_e = 9.1 \times 10^{-31}$ kg, $q = e = 1.6
\times 10^{-19}$ C) at distance $r = 1$ m.

\textbf{GR:} Gravitational deflection angle $\alpha_{\text{GR}} = 4Gm_e/(c^2 r)
\sim 10^{-54}$ rad (unmeasurable).

\textbf{PM:} Total deflection includes electrostatic term:



\[
\alpha_{\text{PM}}
= \alpha_G + \alpha_E
= \frac{\mu_G G m_e}{r} + \frac{\mu_E e}{4\pi\varepsilon_0 r}.
\]



With dimensional estimates, $\alpha_E / \alpha_G \sim 10^{21}$ (charge dominates
vastly over mass for electrons). However, both are still unmeasurably small in
absolute terms.

\textbf{Significance:} PM predicts the \emph{sign and structure} differ for
opposite charges, while GR treats them identically.

\subsection*{2.2 Plasma Lensing by Galaxy Clusters}

\textbf{Scenario:} Gravitational lensing through ionized intracluster medium
with $n_e \sim 10^3$ cm$^{-3}$.

\textbf{GR:} Lensing depends only on cluster mass distribution. Plasma causes
dispersion (group delay) but doesn't modify lensing geometry.

\textbf{PM:} Lensing modified by plasma term:



\[
n_{\text{tot}}
= 1 + \mu_G\,\Phi_G - \alpha_{\text{plasma}}\,n_e.
\]



For path length $L \sim 1$ Mpc through $n_e \sim 10^9$ m$^{-3}$:



\[
\delta n_{\text{plasma}}
\sim -\alpha_{\text{plasma}}\,n_e
\sim -10^{-15}\times 10^9
\sim -10^{-6}.
\]



Comparable to $\delta n_G \sim \mu_G M / r \sim 10^{-6}$ for cluster-scale
lensing.

\textbf{Significance:} Frequency-dependent lensing (not just dispersion) could
be observable in radio vs optical lensing of same source.

\subsection*{2.3 Laboratory Solenoid Test}

\textbf{Scenario:} Light beam passing near 1 MA superconducting coil at 10 cm
distance.

\textbf{GR:} No measurable effect. EM field energy contributes $\sim 10^{-40}$
to gravitational field.

\textbf{PM (naive):} If $\kappa_B \sim c$, predicts $\sim 10$ radian deflection
(absurd).

\textbf{PM (constrained):} Observational null result requires either:

\begin{itemize}
    \item $\kappa_B < 10^{-10}\,c$ (magnetic potential coupling dead), or
    \item $A_q < 10^{-10}$ m$^3$/(C$\cdot$s) (charge current coupling tiny).
\end{itemize}

Resolution: drop $\kappa_B$ term entirely; constrain $A_q$.

\textbf{Significance:} PM's unified structure \emph{predicted} an interaction,
then observations \emph{constrained} it. This is how science should work.

\section*{3. The Observational Hierarchy}

PM reveals a natural hierarchy of medium couplings:



\[
\begin{array}{lcl}
\text{Mass density} & \rightarrow \phi: & \mu_G \sim 10^{-27}\,\text{m}^2/\text{kg} \\
\text{Mass current} & \rightarrow \mathbf u: & A_M \sim 10^{-28}\,\text{m}^3/(\text{kg}\cdot\text{s}) \\
\text{Charge density} & \rightarrow \phi: & \mu_E \sim 10^{-37}\,\text{m}^2/\text{C} \text{ (est.)} \\
\text{Charge current} & \rightarrow \mathbf u: & A_q \lesssim 10^{-10}\,\text{m}^3/(\text{C}\cdot\text{s}) \\
\text{Plasma density} & \rightarrow \phi: & \alpha_{\text{plasma}} \sim 10^{-15}\,\text{m}^3 \text{ (est.)}
\end{array}
\]



\textbf{Pattern:} Mass couples $\sim 10^{17}$--$10^{37}$ times more strongly
than charge (depending on which sector).

This explains why:

\begin{itemize}
    \item Gravitational physics dominates in all tested regimes.
    \item GR (which omits EM-gravity cross-terms) works so well.
    \item PM cross-terms have not been observed yet.
\end{itemize}

But the \emph{structure} exists, making PM richer than GR.

\section*{4. Where PM and GR Make Different Predictions}

\subsection*{4.1 Already Tested (Consistent)}

\begin{itemize}
    \item Weak-field lensing: PM matches GR ($\gamma = \beta = 1$).
    \item Solar system tests: PM reproduces all GR predictions.
    \item Binary pulsar timing: PM frame-dragging matches GR.
\end{itemize}

\subsection*{4.2 Not Yet Tested (PM Predicts Differences)}

\begin{itemize}
    \item \textbf{Strong gravitational fields:} PM has no singularities (index
          diverges, but spacetime remains flat). GR has black hole horizons and
          singularities.
    \item \textbf{Charged-mass lensing:} PM predicts charge-dependent corrections;
          GR predicts charge-independence.
    \item \textbf{Plasma-modified lensing:} PM predicts electron density affects
          deflection angle; GR predicts only dispersion.
    \item \textbf{GW wave speeds:} PM allows $c_\phi \neq c_u$; GR requires $c_{GW} = c$.
    \item \textbf{Current-induced effects:} PM predicts (tiny) dragging from
          extreme currents; GR predicts none beyond stress-energy.
\end{itemize}

\section*{5. Philosophical Implications}

\subsection*{5.1 No Special Cases}

PM treats gravity, EM, and plasma as \emph{different excitation channels} of a
single medium. The equations have the same form; only the coupling strengths
differ. This is analogous to how:

\begin{itemize}
    \item All EM phenomena emerge from Maxwell's equations.
    \item All particle interactions in QFT come from gauge symmetries.
\end{itemize}

\subsection*{5.2 Emergence vs Fundamentality}

\begin{itemize}
    \item \textbf{GR view:} Spacetime geometry is fundamental; matter couples to it.
    \item \textbf{PM view:} Flat spacetime + medium is fundamental; geometry emerges
          from optical properties.
\end{itemize}

Both give the same predictions in weak fields, but suggest different physics in
extreme regimes.

\subsection*{5.3 Unity Through Medium Properties}

PM suggests the apparent separation of forces is an \emph{accident of coupling
strengths}, not a fundamental distinction. If $A_q \sim A_M$, gravitational and
electromagnetic dragging would be equally important—we'd have noticed the
unified structure long ago.

Instead, nature chose $A_q \ll A_M$, hiding the unity and making forces
\emph{appear} separate. PM reveals the hidden structure.

\section*{6. Conclusion}

The Pushing--Medium theory predicts cross-force interactions that General
Relativity's geometric-vs-gauge structure cannot accommodate:

\begin{itemize}
    \item Unified index contributions from mass, charge, and plasma.
    \item Current-induced dragging from mass and charge flows.
    \item Dynamic field-level energy exchange via compressibility.
    \item Multiple gravitational wave speeds.
\end{itemize}

Observations constrain these interactions to be negligible in all tested
regimes, explaining GR's success. But the interactions exist \emph{in
principle}, opening windows for precision tests and revealing a potential
underlying unity that GR's structure cannot express.

PM demonstrates that a ``no special cases'' philosophy—treating all forces as
medium excitations with empirically determined coupling strengths—produces a
mathematically consistent, observationally viable, and conceptually unified
framework that extends beyond GR while remaining compatible with it in the
weak-field regime.

\end{document}
